\chapter{Les alcools~: activer le groupe OH}\label{ch:b1:alcohol-activation}

Les alcools sont partout~: produits par tonnes à partir des alcènes et
par fermentation, fournis par chaque addition de Grignard du chapitre
précédent. Pourtant, un alcool ne subit pas de substitution comme le
fait un halogénoalcane~: le groupe OH devrait partir sous forme d'ion
hydroxyde, une \omterm{def:b1:predominant-reaction:strong-acid}{base forte} et un très mauvais \omterm{def:b1:electronic-effects:nucleophile}{groupe partant}
(\cref{ch:b1:nucleophilic-substitution}). Les chimistes ont trois
réponses. Retirer plutôt le proton du OH~: l'\omterm{def:b1:alcohol-activation:alkoxide}{alcoolate} devient un
\omterm{def:b1:electronic-effects:nucleophile}{nucléophile} puissant. Ajouter un proton~: le \omterm{def:b1:electronic-effects:nucleophile}{groupe partant} devient
l'eau. Ou remplacer l'hydrogène par un groupe qui intègre l'oxygène à un
excellent \omterm{def:b1:electronic-effects:nucleophile}{groupe partant}, sans toucher au carbone. Ce chapitre suit ces
trois voies, ainsi que les choix de mécanisme et de stéréochimie qui les
accompagnent.

\begin{recall}
SN1, SN2 et leur stéréochimie (\cref{ch:b1:nucleophilic-substitution})~;
E1, E2 et la \omterm{prop:b1:elimination:zaitsev}{règle de Zaïtsev} (\cref{ch:b1:elimination})~; les
\omterm{def:b1:predominant-reaction:acidity-constant}{constantes d'acidité} et la méthode de la \omterm{def:b1:predominant-reaction:predominant-reaction}{réaction prépondérante}
(\cref{ch:b1:predominant-reaction})~; l'effet des groupes alkyle sur
l'acidité des alcools (\cref{exo:b1:electronic-effects:11}).
\end{recall}

\section{Les alcools en tant qu'acides~: les alcoolates}

\begin{definition}[Alcoolate]\label{def:b1:alcohol-activation:alkoxide}
Un \emph{alcoolate}\index{alcoolate} est la base conjuguée \ce{RO-} d'un
alcool \ce{ROH}~: méthanolate \ce{CH3O-}, éthanolate \ce{CH3CH2O-},
\textit{tert}-butanolate \ce{(CH3)3CO-}. Les alcoolates sont des \omterm{def:b1:predominant-reaction:strong-acid}{bases
fortes} et de bons \omterm{def:b1:electronic-effects:nucleophile}{nucléophiles}.
\end{definition}

\begin{proposition}[Acidité des alcools]\label{prop:b1:alcohol-activation:alcohol-acidity}
Les alcools sont des acides plus faibles que l'eau, ou à peu près aussi
faibles~: $\mathrm{p}K_a$ de
15.5 (méthanol), 15.9 (éthanol), 17.1 (propan-2-ol), 19.2
(2-méthylpropan-2-ol), contre 14.0 pour le couple \ce{H2O}/\ce{OH-}.
L'hydroxyde ne les déprotone donc que partiellement~; une conversion
complète en \omterm{def:b1:alcohol-activation:alkoxide}{alcoolate} exige une base plus forte ou un métal alcalin.
\end{proposition}
% ledger: pka:methanol, pka:ethanol, pka:2-propanol, pka:tBuOH

\begin{proof}
La réaction \ce{ROH + OH- <=> RO- + H2O} a pour constante $K = 10^{14.0 -
\mathrm{p}K_a(\ce{ROH})}$, comprise entre $10^{-1.5}$ et $10^{-5}$~:
elle est défavorable. Le sodium métallique réduit le proton,
\ce{2ROH + 2Na -> 2RO- + 2Na+ + H2}, et l'ion hydrure de l'hydrure de
sodium le capte de façon irréversible,
\ce{ROH + H- -> RO- + H2}~: le dihydrogène s'échappe, et la réaction est
totale quelle que soit sa constante d'équilibre.
\end{proof}

\begin{inthelab}[L'hydrure de sodium]
L'hydrure de sodium est vendu sous forme de poudre grise dispersée dans
une huile minérale, qui le protège de l'humidité de l'air. Il réagit
violemment avec l'eau en libérant du dihydrogène inflammable~; on le
pèse rapidement, on l'ajoute par portions à l'alcool anhydre sous gaz
inerte, et le dégagement de dihydrogène montre l'\omterm{def:b1:extent-q-and-k:extent}{avancement} de la
déprotonation. L'excès d'hydrure est détruit à la fin par addition lente
d'un alcool, jamais d'eau.
\end{inthelab}

\section{La synthèse de Williamson}

\begin{definition}[Synthèse de Williamson]\label{def:b1:alcohol-activation:williamson}
La \emph{synthèse de Williamson}\index{synthese de Williamson@synthèse de Williamson}
prépare un éther $\mathrm{R{-}O{-}R'}$ par une réaction SN2 d'un
\omterm{def:b1:alcohol-activation:alkoxide}{alcoolate} \ce{RO-} sur un halogénoalcane (ou un \omterm{def:b1:alcohol-activation:sulfonate}{ester sulfonique})
$\mathrm{R'{-}X}$.
\end{definition}

\begin{omfigure}
\setchemfig{atom sep=2.1em}
\schemestart
\chemfig{CH_3CH_2-@{o}O^{-}}\,\chemfig{Na^{+}}
\+
\chemfig{@{c}CH_3-[@{ci}]@{i}I}
\arrow{->}
\chemfig{CH_3CH_2-O-CH_3}
\+
\chemfig{Na^{+}}\,\chemfig{I^{-}}
\schemestop

\medskip
\chemmove{\draw[->, thick, shorten <=2pt, shorten >=2pt](o) .. controls +(80:8mm) and +(100:8mm) .. (c);
          \draw[->, thick, shorten <=2pt, shorten >=2pt](ci) .. controls +(90:5mm) and +(90:5mm) .. (i);}

{\small \omterm{def:b1:alcohol-activation:williamson}{Synthèse de Williamson} du méthoxyéthane~: l'éthanolate attaque le
carbone de l'iodométhane du côté opposé à l'iode (SN2).}
\end{omfigure}

\begin{method}[Choisir les deux partenaires]\label{met:b1:alcohol-activation:williamson-choice}
Un éther $\mathrm{R{-}O{-}R'}$ peut être coupé de part et d'autre de
l'oxygène.
\begin{enumerate}
  \item Écrire les deux couples possibles~: \ce{RO-} avec
        $\mathrm{R'X}$, et
        $\mathrm{R'O^-}$ avec \ce{RX}.
  \item Garder le couple dans lequel l'halogénure est méthylique ou
        primaire~: la SN2 exige un carbone peu encombré.
  \item Rejeter un couple comportant un halogénure tertiaire~:
        l'\omterm{def:b1:alcohol-activation:alkoxide}{alcoolate}, \omterm{def:b1:predominant-reaction:strong-acid}{base forte}, provoquerait plutôt une élimination (E2).
  \item Avec un halogénure secondaire, s'attendre à un mélange d'éther et
        d'alcène.
\end{enumerate}
\end{method}

\begin{example}[Un éther dissymétrique]\label{ex:b1:alcohol-activation:williamson}
Le 2-méthoxy-2-méthylpropane (méthyl \textit{tert}-butyl éther)~: le
\textit{tert}-butanolate de sodium avec l'iodométhane le donne
proprement~; le méthanolate de sodium avec le bromure tertiaire, le
2-bromo-2-méthyl\-propane, donne du 2-méthylpropène par E2.
L'encombrement du côté de l'\omterm{def:b1:alcohol-activation:alkoxide}{alcoolate}, la petite taille du côté de
l'halogénure.
\end{example}

\section{Activer le groupe OH}

\begin{definition}[Esters sulfoniques]\label{def:b1:alcohol-activation:sulfonate}
Un \emph{ester sulfonique}\index{ester sulfonique} $\mathrm{R{-}O{-}SO_2{-}R'}$ se forme
à partir d'un alcool et d'un chlorure de sulfonyle $\mathrm{R'SO_2Cl}$
en présence d'une base (la pyridine). Le \emph{tosylate}\index{tosylate},
\ce{R-OTs}, provient du chlorure de 4-méthylbenzènesulfonyle (chlorure de
tosyle)~; le \emph{mésylate}\index{mesylate@mésylate}, \ce{R-OMs}, du chlorure de
méthanesulfonyle. L'anion sulfonate $\mathrm{R'SO_3^-}$, base conjuguée
d'un \omterm{def:b1:predominant-reaction:strong-acid}{acide fort} et stabilisé par mésomérie sur trois oxygènes, est un
excellent \omterm{def:b1:electronic-effects:nucleophile}{groupe partant}.
\end{definition}

\begin{omfigure}
\setchemfig{atom sep=2em}
\schemestart
\chemfig{R-O-H}
\+
\chemfig{Cl-S(=[2]O)(=[6]O)-[:0]*6(-=-(-CH_3)=-=)}
\arrow{->[pyridine]}
\chemfig{R-O-S(=[2]O)(=[6]O)-[:0]*6(-=-(-CH_3)=-=)}
\schemestop
\par\vspace{6pt}

{\small Tosylation d'un alcool. La liaison O--H est remplacée par O--S~;
la liaison C--O de l'alcool n'est pas touchée. La pyridine capte le HCl
formé.}
\end{omfigure}

\begin{proposition}[La configuration au cours de la tosylation et de la substitution]\label{prop:b1:alcohol-activation:tosylation-retention}
La tosylation d'un alcool sur un carbone stéréogène conserve sa
\omterm{def:b1:stereochemistry-in-depth:isomers}{configuration} (rétention)~; une substitution SN2 ultérieure du \omterm{def:b1:alcohol-activation:sulfonate}{tosylate}
l'inverse. Les deux étapes réunies remplacent OH par un \omterm{def:b1:electronic-effects:nucleophile}{nucléophile} avec
inversion.
\end{proposition}

\begin{proof}
Au cours de la tosylation, les liaisons qui changent sont O--H et
S--Cl~: l'oxygène de l'alcool attaque le soufre et son hydrogène passe à
la base. Aucune liaison au carbone stéréogène n'est rompue, si bien que
la disposition autour de lui reste la même. Lors de l'étape SN2, la
liaison C--O se rompt pendant que le \omterm{def:b1:electronic-effects:nucleophile}{nucléophile} se lie du côté opposé
(\cref{prop:b1:nucleophilic-substitution:sn2-inversion}).
\end{proof}

\begin{omfigure}
\begin{tikzpicture}[font=\small, >={Stealth}]
  \node (a) at (0,0) {\chemfig{C(-[:150]H_3C)(<[:210]H)(<:[:250]C_2H_5)-OH}};
  \node at (0,-1.4) {\cip{S}-butan-2-ol};
  \draw[->, thick] (1.3,0) -- (2.5,0) node[midway, above] {TsCl} node[midway, below, font=\scriptsize] {rétention};
  \node (b) at (3.9,0) {\chemfig{C(-[:150]H_3C)(<[:210]H)(<:[:250]C_2H_5)-OTs}};
  \node at (3.9,-1.4) {tosylate \cip{S}};
  \draw[->, thick] (5.3,0) -- (6.6,0) node[midway, above] {\ce{CN-}} node[midway, below, font=\scriptsize] {SN2, inversion};
  \node (c) at (8.3,0) {\chemfig{[:0]NC-C(-[:30]CH_3)(<[:-30]H)(<:[:-70]C_2H_5)}};
  \node at (8.3,-1.4) {nitrile \cip{R}};
\end{tikzpicture}

{\small Du \cip{S}-butan-2-ol au 2-méthylbutanenitrile~: la tosylation
conserve la \omterm{def:b1:stereochemistry-in-depth:isomers}{configuration}, l'étape SN2 avec le cyanure l'inverse.
Bilan~: OH remplacé par CN avec inversion.}
\end{omfigure}

Le même objectif peut être atteint en milieu acide~: le groupe OH,
protoné, part sous forme d'eau, \ce{R-OH2+ -> R+ + H2O} (SN1, E1) ou
sous l'effet d'une attaque (SN2). Mais l'acide restreint le choix des
\omterm{def:b1:electronic-effects:nucleophile}{nucléophiles} à ceux qui lui résistent --- ions halogénure, eau, alcools
--- et ouvre la porte aux \omterm{def:b1:electronic-effects:intermediates}{carbocations} et à leurs réactions secondaires.

\section{Conversion en halogénoalcanes}

\begin{proposition}[Alcools et halogénures d'hydrogène]\label{prop:b1:alcohol-activation:hx-by-class}
Les alcools tertiaires donnent l'halogénoalcane avec l'acide
chlorhydrique concentré à température ambiante (SN1)~; les alcools
primaires exigent l'acide bromhydrique ou iodhydrique et un chauffage
(SN2 sur l'alcool protoné)~; les alcools secondaires réagissent à des
vitesses intermédiaires, souvent par les deux mécanismes.
\end{proposition}

\begin{proof}
La protonation transforme OH en \ce{-OH2+}. Pour un alcool tertiaire,
l'eau part facilement, ce qui donne un \omterm{def:b1:electronic-effects:intermediates}{carbocation} stable capturé par
l'halogénure (SN1)~: même le modeste \omterm{def:b1:electronic-effects:nucleophile}{nucléophile} \ce{Cl-} suffit. Un
\omterm{def:b1:electronic-effects:intermediates}{carbocation} primaire est trop instable~; l'halogénure doit chasser l'eau
par l'arrière (SN2), ce qui exige un bon \omterm{def:b1:electronic-effects:nucleophile}{nucléophile} (\ce{Br-}, \ce{I-},
meilleurs que \ce{Cl-} en milieu protique,
\cref{prop:b1:electronic-effects:nucleophilicity-basicity}) et un
chauffage.
\end{proof}

\begin{example}[Un test de classe]\label{ex:b1:alcohol-activation:lucas}
Un mélange d'acide chlorhydrique concentré et de chlorure de zinc (un
\omterm{def:b1:lewis-resonance-vsepr:lewis-acid}{acide de Lewis} qui aide le OH à partir) trouble aussitôt un alcool
tertiaire --- le chlorure, insoluble, se sépare ---, un alcool secondaire
en quelques minutes, et un alcool primaire presque pas à température
ambiante~: l'ordre de stabilité des \omterm{def:b1:electronic-effects:intermediates}{carbocations} rendu visible dans un
tube à essais.
\end{example}

\section{Déshydratation et formation d'éthers en milieu acide}

\begin{definition}[Déshydratation d'un alcool]\label{def:b1:alcohol-activation:dehydration}
La \emph{déshydratation}\index{deshydratation@déshydratation} d'un alcool est l'élimination
d'eau, catalysée par un \omterm{def:b1:predominant-reaction:strong-acid}{acide fort} (acide sulfurique ou phosphorique) à
chaud, qui donne un alcène~:
$\mathrm{R_2CH{-}CR'_2OH} \to \mathrm{R_2C{=}CR'_2} + \ce{H2O}$.
\end{definition}

\begin{omfigure}
\setchemfig{atom sep=2.1em}
\schemestart
\chemfig{C(-[2]CH_3)(-[:210]H_3C)(-[:-30]CH_2CH_3)-OH}
\arrow{->[\ce{H+}]}
\chemfig{C(-[2]CH_3)(-[:210]H_3C)(-[:-30]CH_2CH_3)-OH_2^{+}}
\arrow{->[\ce{-H2O}]}
\chemfig{C^{+}(-[2]CH_3)(-[:210]H_3C)(-[:-30]CH_2CH_3)}
\schemestop

\medskip
\schemestart
\chemfig{C^{+}(-[2]CH_3)(-[:210]H_3C)(-[:-30]CH_2CH_3)}
\arrow{->[\ce{-H+}][(majoritaire)]}
\chemfig{C(-[2]CH_3)(-[:210]H_3C)=[:-30]CHCH_3}
\schemestop

\medskip
{\small \omterm{def:b1:alcohol-activation:dehydration}{Déshydratation} acido-catalysée du 2-méthylbutan-2-ol (E1)~:
protonation, perte d'eau conduisant au \omterm{def:b1:electronic-effects:intermediates}{carbocation} tertiaire, perte d'un
proton. Arracher le proton du \ce{CH2} donne le 2-méthylbut-2-ène,
trisubstitué (Zaïtsev, majoritaire)~; l'arracher à un groupe méthyle
donne le 2-méthylbut-1-ène (minoritaire).}
\end{omfigure}

\begin{proposition}[Éther ou alcène]\label{prop:b1:alcohol-activation:ether-vs-alkene}
Un alcool primaire chauffé avec l'acide sulfurique donne surtout un
éther symétrique à température modérée (une molécule en substituant une
autre par SN2) et surtout un alcène à plus haute température~; les
alcools secondaires et tertiaires donnent des alcènes.
\end{proposition}

\begin{proof}
Pour un alcool primaire, l'alcool protoné peut être attaqué soit sur le
carbone par une seconde molécule d'alcool (SN2, qui donne \ce{R-O-R}
après perte d'un proton), soit sur un hydrogène en $\beta$ (E2).
L'élimination produit plus de molécules qu'elle n'en consomme et est
favorisée quand la température augmente
(\cref{prop:b1:elimination:sn-vs-e})~; la substitution domine lorsque la
température est plus basse. Les alcools secondaires et tertiaires
s'ionisent en \omterm{def:b1:electronic-effects:intermediates}{carbocations}, qui perdent un proton plus facilement qu'ils
ne sont capturés par un \omterm{def:b1:electronic-effects:nucleophile}{nucléophile} faible, d'autant plus qu'on chauffe~;
l'eau formée est aussi éliminée par distillation de l'alcène volatil.
\end{proof}

\section{Exercices}

\begin{exercise}[$\star$]\label{exo:b1:alcohol-activation:1}
Écrire les réactions de l'éthanol avec le sodium, avec l'hydrure de
sodium et avec l'hydroxyde de sodium~; calculer la constante de la
dernière ($\mathrm{p}K_a$ de l'éthanol~: 15.9).
\end{exercise}

\begin{exercise}[$\star$]\label{exo:b1:alcohol-activation:2}
Donner les produits~: du méthanolate de sodium avec le 1-bromopropane~;
de l'éthanolate de sodium avec l'iodométhane~; du phénolate de sodium
avec le bromoéthane.
\end{exercise}

\begin{exercise}[$\star$]\label{exo:b1:alcohol-activation:3}
Écrire la tosylation du propan-1-ol et la réaction du \omterm{def:b1:alcohol-activation:sulfonate}{tosylate} avec
l'iodure de sodium. Quel est l'avantage par rapport à la réaction
directe de l'alcool avec l'iodure de sodium~?
\end{exercise}

\begin{exercise}[$\star$]\label{exo:b1:alcohol-activation:4}
Donner les alcènes formés par \omterm{def:b1:alcohol-activation:dehydration}{déshydratation} acide du butan-2-ol et du
2-méthylpentan-2-ol, ainsi que l'alcène majoritaire.
\end{exercise}

\begin{exercise}[$\star\star$]\label{exo:b1:alcohol-activation:5}
Proposer une \omterm{def:b1:alcohol-activation:williamson}{synthèse de Williamson} du 1-éthoxy-2-méthylpropane et
justifier le choix des partenaires. Pourquoi l'autre choix donnerait-il
un alcène~?
\end{exercise}

\begin{exercise}[$\star\star$]\label{exo:b1:alcohol-activation:6}
Le \cip{R}-octan-2-ol est converti en son \omterm{def:b1:alcohol-activation:sulfonate}{tosylate}, puis traité par
l'éthanolate de sodium dans l'éthanol. Donner la \omterm{def:b1:stereochemistry-in-depth:isomers}{configuration} de
l'éther. Qu'obtiendrait-on en chauffant le même alcool avec de l'acide
sulfurique~?
\end{exercise}

\begin{exercise}[$\star\star$]\label{exo:b1:alcohol-activation:7}
Écrire le mécanisme de la réaction du 2-méthylpropan-2-ol avec l'acide
chlorhydrique concentré et expliquer pourquoi la réaction est rapide à
température ambiante.
\end{exercise}

\begin{exercise}[$\star\star$]\label{exo:b1:alcohol-activation:8}
Écrire le mécanisme de la formation de l'éthoxyéthane à partir de
l'éthanol dans l'acide sulfurique, ainsi que l'élimination concurrente.
\end{exercise}

\begin{exercise}[$\star\star$]\label{exo:b1:alcohol-activation:9}
Dans le test de classe de ce chapitre, quel alcool, parmi le
butan-1-ol, le butan-2-ol et le 2-méthylpropan-2-ol, se trouble le
premier~? Écrire le produit.
\end{exercise}

\begin{exercise}[$\star\star\star$]\label{exo:b1:alcohol-activation:10}
Le 3,3-diméthylbutan-2-ol, chauffé en milieu acide, donne surtout du
2,3-diméthylbut-2-ène, dont le squelette carboné diffère de celui de
l'alcool. Proposer un mécanisme comportant une migration-1,2 d'un groupe
méthyle dans le \omterm{def:b1:electronic-effects:intermediates}{carbocation}, et justifier cette migration.
\end{exercise}

\begin{exercise}[$\star\star\star$]\label{exo:b1:alcohol-activation:11}
Concevoir une synthèse du \cip{S}-2-méthoxybutane à partir du
\cip{R}-butan-2-ol, et une autre à partir du \cip{S}-butan-2-ol, chacune
gardant la \omterm{def:b1:stereochemistry-in-depth:isomers}{configuration} sous contrôle.
\end{exercise}

\begin{exercise}[$\star\star\star$]\label{exo:b1:alcohol-activation:12}
Montrer, à l'aide du \omterm{def:b1:extent-q-and-k:quotient}{quotient de réaction} (\cref{ch:b1:extent-q-and-k}),
que l'équilibre \ce{2CH3CH2OH <=> (CH3CH2)2O + H2O} peut être déplacé
vers l'éther en éliminant l'eau en continu.
\end{exercise}

\section{Problème~: deux voies vers un éther}

\begin{problem}[{Problème du week-end --- les deux voies de Williamson
vers le méthyl \textit{tert}-butyl éther, un éther à stéréochimie
contrôlée issu d'un alcool optiquement actif, la \omterm{def:b1:alcohol-activation:dehydration}{déshydratation} d'un
alcool tertiaire, et la masse d'éther obtenue}]
\label{pb:b1:alcohol-activation:1}
Le 2-méthoxy-2-méthylpropane (méthyl \textit{tert}-butyl éther, MTBE) a
été produit à grande échelle comme additif des carburants. Masses
molaires (\unit{g/mol})~: \ce{C4H10O}
74.0, \ce{C5H12O} 88.0, \ce{CH3I} 141.9.

\medskip
\noindent\textbf{Partie I --- Le MTBE par la \omterm{def:b1:alcohol-activation:williamson}{synthèse de Williamson}.}
\begin{enumerate}
  \item Écrire les deux couples de partenaires (\omterm{def:b1:alcohol-activation:alkoxide}{alcoolate} et halogénure)
        qui pourraient donner le MTBE.
  \item Quel couple échoue~? Écrire la réaction qui a lieu à la place.
  \item Écrire la réaction qui réussit et son mécanisme.
  \item Comment prépare-t-on le \textit{tert}-butanolate de sodium à
        partir du 2-méthylpropan-2-ol~? Pourquoi pas avec l'hydroxyde de
        sodium~?
  \item Pourquoi mène-t-on la réaction dans l'alcool anhydre ou dans un
        \omterm{def:b1:intermolecular-forces-solvents:solvent-classes}{solvant aprotique}~?
  \item Le produit est-il \omterm{def:b1:stereochemistry-in-depth:stereocentre}{chiral}~?
\end{enumerate}

\medskip
\noindent\textbf{Partie II --- Un éther optiquement actif.}
\begin{enumerate}[resume]
  \item Le \cip{R}-butan-2-ol est converti en son \omterm{def:b1:alcohol-activation:sulfonate}{tosylate}. Écrire la
        réaction et donner la \omterm{def:b1:stereochemistry-in-depth:isomers}{configuration} du \omterm{def:b1:alcohol-activation:sulfonate}{tosylate}.
  \item Le \omterm{def:b1:alcohol-activation:sulfonate}{tosylate} réagit avec le méthanolate de sodium. Écrire le
        produit et sa \omterm{def:b1:stereochemistry-in-depth:isomers}{configuration}.
  \item Quelle \omterm{def:b1:stereochemistry-in-depth:isomers}{configuration} de l'alcool donne le
        \cip{R}-2-méthoxybutane par la même voie~?
  \item Autre voie~: l'\omterm{def:b1:alcohol-activation:alkoxide}{alcoolate} du \cip{R}-butan-2-ol avec
        l'iodométhane. \omterm{def:b1:stereochemistry-in-depth:isomers}{Configuration} du produit~? Pourquoi~?
  \item Quelle réaction secondaire est en concurrence à la question 8,
        et pourquoi est-elle limitée ici~?
  \item Pourquoi le chauffage du \cip{R}-butan-2-ol avec du méthanol dans
        l'acide sulfurique donnerait-il un éther presque racémique, s'il
        en donne~?
\end{enumerate}

\medskip
\noindent\textbf{Partie III --- La \omterm{def:b1:alcohol-activation:dehydration}{déshydratation} d'un alcool tertiaire.}
\begin{enumerate}[resume]
  \item Écrire le mécanisme de la \omterm{def:b1:alcohol-activation:dehydration}{déshydratation} acide du
        2-méthylbutan-2-ol.
  \item Donner les deux alcènes et prévoir l'alcène majoritaire.
  \item Pourquoi aucun éther ne se forme-t-il à partir de cet alcool~?
  \item Pourquoi distille-t-on le 2-méthylbut-2-ène au fur et à mesure de
        sa formation~?
  \item Tracer le \omterm{def:b1:elementary-steps:energy-profile}{profil énergétique} de l'étape E1 et y repérer l'\omterm{def:b1:elementary-steps:rds}{étape
        cinétiquement déterminante}.
  \item Un alcool primaire pourrait-il être déshydraté par le même
        mécanisme~?
\end{enumerate}

\medskip
\noindent\textbf{Partie IV --- Les masses.}
\begin{enumerate}[resume]
  \item On convertit \qty{10.0}{g} de 2-méthylpropan-2-ol en \omterm{def:b1:alcohol-activation:alkoxide}{alcoolate}.
        Calculer la quantité de matière.
  \item Quelle masse d'iodométhane faut-il pour un excès de
        \qty{10}{\percent}~?
  \item Calculer la masse théorique de MTBE.
  \item Calculer la masse obtenue pour un \omterm{def:b1:extent-q-and-k:fractional-extent}{rendement} de \qty{75}{\percent}.
\end{enumerate}
\end{problem}
